\section{Metodología}
\label{sec:metodologia}

La verificación de un sistema en tiempo real exige métodos distintos a los bancos de
consultas de los reportes anteriores: aquí se combinan pruebas inyectadas
reproducibles (sin micrófono ni cámara), pruebas físicas con estímulos controlados y
mediciones de latencia sobre el hardware real.

\subsection{Entorno}

Todas las mediciones se realizaron en el equipo de desarrollo del proyecto: GPU
NVIDIA RTX 2050 (usada exclusivamente por Whisper, con carga explícita de las
bibliotecas CUDA de distribución pip y validación por inferencia de silencio al
arrancar, con respaldo a CPU int8 si la GPU no responde) y CPU para el encoder
ResNet18 (suficiente a $128\times128$) y todas las operaciones de memoria.
Los modelos del experimento base se cargan una única vez ($\sim$30 s) y permanecen en
solo lectura.

\subsection{Pruebas inyectadas (sin hardware)}

\begin{description}
    \item[Inyección de audio (\texttt{--wavtest}).] Un archivo WAV se alimenta por el
    mismo camino del micrófono (remuestreo a 16 kHz $\to$ VAD $\to$ transcripción
    $\to$ enrutamiento $\to$ evocación), lo que permite medir la cadena completa de
    voz de forma reproducible y sin variabilidad del hablante.

    \item[Autoprueba sin ventana (\texttt{--selftest}).] El sistema corre $n$ segundos
    sin interfaz, reporta fps, RMS inicial del micrófono y estado del canal visual, y
    guarda un cuadro anotado con todos los overlays para inspección.

    \item[Cuadros sintéticos.] Imágenes ETH-80 reales colocadas digitalmente dentro
    del recuadro de análisis (llenándolo al 55\%) y cuadros vacíos, para validar por
    separado el enrutamiento visual y el rechazo sin la variabilidad de la cámara.

    \item[Detección de bloqueos del sistema operativo.] Sondas al arrancar: RMS del
    micrófono $\approx 0$ dispara una advertencia con instrucciones (el permiso de
    privacidad de Windows para aplicaciones de escritorio aplica también a WASAPI); el
    bloqueo de cámara entrega un cuadro gris con candado que el directorio rechaza
    ---comportamiento correcto que la autoprueba hace visible.
\end{description}

\subsection{Pruebas físicas de percepción}
\label{sec:met_fisica}

Se probaron dos formas de presentar estímulos ETH-80 a la cámara web:

\begin{enumerate}
    \item \textbf{Pantalla de teléfono} mostrando imágenes pre-validadas
    (160 imágenes de $512^2$).
    Falló de forma consistente: brillo de pantalla, moiré y encuadre colapsan el
    enrutamiento.
    El fallo fue \emph{anticipado} por una simulación previa que degradaba
    sistemáticamente encuadre y brillo, y una segunda simulación midió qué rescates
    eran posibles.

    \item \textbf{Hoja impresa} (4 páginas carta a 300 dpi, 16 imágenes pre-validadas
    a $\sim$8.5 cm por lado, papel mate).
    El papel elimina brillo y moiré; con el recuadro ajustable la impresión se llena a
    distancia de enfoque cómoda.
    El diagnóstico de la primera sesión física (solo \textit{cow} fiable) condujo a la
    política de color de la Sección~\ref{sec:arquitectura}.
\end{enumerate}

\subsubsection*{Banco de validación de la política de color}

La corrección~\eqref{eq:colorfix} se validó fuera de línea sobre 24 imágenes ETH-80
(3 por clase) bajo tres condiciones simuladas de cámara: limpia, degradación grisácea
media y fuerte (contracción de contraste y corrimiento de canal medidos en la cámara
real).
Se reporta el enrutamiento correcto con la política completa (crudo, y rescate solo
tras rechazo) contra la línea base sin rescate.

\subsection{Caracterización de los métodos de pista faltante}
\label{sec:met_pista}

Los métodos RS, ST y SS y el prototipo se caracterizaron sobre la misma proyección por
pista, en dos escenarios:

\begin{enumerate}
    \item \textbf{Fuera de línea} (prueba de humo reproducible): pistas de un token
    (\texttt{pear} $\to$ agente \textit{pear}; \texttt{car} $\to$ agente \textit{car})
    con los cuatro valores de distancia~\eqref{eq:test}, el peso medio del patrón y el
    tiempo de cómputo por método.

    \item \textbf{En las interfaces}: la aplicación de trazado (consulta
    ``a ripe green pear from the orchard'', suite completa sobre el token ganador) y
    la ventana en tiempo real (frase ``animal with a mane'' $\to$ \textit{horse},
    prototipo y objeto definido con sus distancias en el panel).
\end{enumerate}

Al ser RS y ST muestreos estocásticos, sus distancias varían entre ejecuciones; se
reportan corridas individuales representativas ---señaladas como tales--- y la
jerarquía cualitativa, que fue estable en todas las ejecuciones observadas.

\subsection{Métricas}

\begin{itemize}
    \item \textbf{Latencia de voz}: transcripción (fin de habla $\to$ texto) y
    enrutamiento (texto $\to$ decisión), por separado y extremo a extremo.
    \item \textbf{Latencia de visión}: tiempo por análisis de cuadro y fps de la
    interfaz con la cámara activa.
    \item \textbf{Latencia de evocación}: fin del enrutamiento $\to$ imagen
    reconstruida (incluyendo prototipo y distancias).
    \item \textbf{Enrutamiento correcto} bajo las condiciones del banco de color
    (aciertos sobre 24).
    \item \textbf{Distancia de retro-proyección}~\eqref{eq:test} como métrica común de
    los métodos de recuperación.
\end{itemize}
